\emph{Preliminaries.} Write $u = h(k)$, $\tau = \tr h$, $Q = Q(k)$. The stored proof of Theorem \claimref{synthesis.gravitational-pair-production-from-the-operator-source-at-first-order}(i) gives, for every $k$ and every symmetric $h$,
\[
\mathrm{Ric}(a) = -\tfrac12\Bigl[(a \cdot k)(u - \tau k) - Q\,h(a) + (a \cdot u)\,k\Bigr], \qquad \tr\mathrm{Ric} = -(k \cdot u - Q\tau),
\]
using $k \cdot h(a) = a \cdot u$. Hence $b \cdot \mathrm{Ric}(a)$ is symmetric in $a, b$, and so is $\mathcal{G}_{k}[h] = \mathrm{Ric} - \tfrac12\,\id\,\tr\mathrm{Ric}$: the Einstein map takes symmetric maps to symmetric maps.

\emph{(i).} Let $Q = 0$. Then $\mathcal{G}_{k}[h](a) = -\tfrac12\bigl[(a \cdot k)(u - \tau k) + (a \cdot u)k\bigr] + \tfrac12 (k \cdot u)\,a$. If this vanishes for all $a$ then $(k \cdot u)\,a$ equals a map of rank at most two for all $a$, so $k \cdot u = 0$, and $(a \cdot k)(u - \tau k) + (a \cdot u)k = 0$ for all $a$. If $u$ were not a multiple of $k$, some $a$ would have $a \cdot k = 1$, $a \cdot u = 0$, giving $u = \tau k$, a contradiction; so $u = ck$, and then $(2c - \tau)(a \cdot k)k = 0$ gives $c = \tfrac12\tau$. Conversely $u = \tfrac12\tau k$ gives $k \cdot u = 0$ and $\mathcal{G}_{k}[h] = 0$. So $\ker\mathcal{G}_{k} = \{h : h(k) = \tfrac12\tr(h)k\}$. This is cut out by four linear conditions on ten dimensions, and they are independent: with $n$ null, $n \cdot k = 1$, $n \cdot m_{i} = 0$, the map $h \mapsto h(k) - \tfrac12\tr(h)k$ sends $\id$ to $-k$, the map $a \mapsto (a \cdot n)n$ to $n$, and $a \mapsto (a \cdot n)m_{i} + (a \cdot m_{i})n$ to $m_{i}$, a basis. So the kernel has dimension six (the rank four of the conditions is also computed). The displacements lie in it (dimension four, Theorem \claimref{synthesis.weyl-flat-variations-carry-only-the-stress-trace-and-the-whole-stress-clause-is-the-trace-free-channel}(i)); $n_{+}$ and $n_{\times}$ have $n(k) = 0$ and trace $m_{1} \cdot m_{1} - m_{2} \cdot m_{2} = 0$ and $2\,m_{1} \cdot m_{2} = 0$, so they lie in it; and $N_{k} \cap G_{k} = 0$: every element of $N_{k}$ kills $n$, while a displacement that kills $n$ has $\xi + (n \cdot \xi)k = 0$, hence $n \cdot \xi = 0$ on pairing with $n$, hence $\xi = 0$. Hence $\ker\mathcal{G}_{k} = G_{k} \oplus N_{k}$. Replacing $m_{i}$ by $m_{i} + c_{i}k$ changes $n_{+}$ and $n_{\times}$ by maps of the form $a \mapsto (a \cdot k)\xi + (a \cdot \xi)k$, and a rotation of $m_{1}, m_{2}$ rotates $n_{+}, n_{\times}$ by twice the angle: with $m_{1}' = c\,m_{1} + s\,m_{2}$, $m_{2}' = -s\,m_{1} + c\,m_{2}$, $c^{2} + s^{2} = 1$, expanding the two products gives $n_{+}' = (c^{2} - s^{2})\,n_{+} + 2cs\,n_{\times}$ and $n_{\times}' = -2cs\,n_{+} + (c^{2} - s^{2})\,n_{\times}$. The norms: $\ip{n_{+}}{n_{+}} = (m_{1} \cdot m_{1})^{2} + (m_{2} \cdot m_{2})^{2} = 2$, likewise for $n_{\times}$.

\emph{(ii).} From the Preliminaries, for symmetric $h, h'$ with $u' = h'(k)$, $\tau' = \tr h'$,
\[
\ip{h'}{\mathcal{G}_{k}[h]} = -u \cdot u' + \tfrac12\tau\,(k \cdot u') + \tfrac12\tau'\,(k \cdot u) + \tfrac12 Q\,\ip{h'}{h} - \tfrac12 Q\,\tau\tau' ,
\]
symmetric in $h, h'$ (the terms: $\ip{h'}{\mathrm{Ric}} = -\tfrac12[2u \cdot u' - \tau\,k \cdot u' - Q\ip{h'}{h}]$ and $\ip{h'}{\id} = \tau'$). A symmetric map for a nondegenerate pairing has image equal to the annihilator of its kernel. By clause (iii) of the stored first-order theorem a symmetric map annihilates $G_{k}$ exactly when it is conserved at $k$ ($\ip{\mathcal{T}}{a \mapsto (a \cdot k)\xi + (a \cdot \xi)k} = 2\,\xi \cdot \mathcal{T}(k)$). So at null $k$ the image is the conserved sources that pair to zero with $n_{+}$ and $n_{\times}$; its dimension is $10 - 6 = 4$.

\emph{(iii).} Complete $k, m_{1}, m_{2}$ to a basis by a null $n$ with $n \cdot k = 1$, $n \cdot m_{i} = 0$. Writing $x \odot y$ for the map $a \mapsto (a \cdot x)y + (a \cdot y)x$, the identity is $\id = k \odot n - \tfrac12 m_{1} \odot m_{1} - \tfrac12 m_{2} \odot m_{2}$, and a symmetric map with $\mathcal{T}(k) = 0$ has no term containing $n$: $\mathcal{T} = \alpha\,k \odot k + \sum_{i}\beta_{i}\,k \odot m_{i} + \tfrac12\gamma_{11}\,m_{1} \odot m_{1} + \tfrac12\gamma_{22}\,m_{2} \odot m_{2} + \gamma_{12}\,m_{1} \odot m_{2}$. Its trace is $-(\gamma_{11} + \gamma_{22})$, its pairings are $\ip{\mathcal{T}}{n_{+}} = \gamma_{11} - \gamma_{22}$, $\ip{\mathcal{T}}{n_{\times}} = 2\gamma_{12}$, so $P_{k}\mathcal{T} = \tfrac12(\gamma_{11} - \gamma_{22})\,n_{+} + \gamma_{12}\,n_{\times}$ with $n_{+} = \tfrac12 m_{1} \odot m_{1} - \tfrac12 m_{2} \odot m_{2}$, $n_{\times} = m_{1} \odot m_{2}$. Then
\[
\mathcal{T} - \tfrac12\,\id\,\tr\mathcal{T} - P_{k}\mathcal{T} = \alpha\,k \odot k + \sum_{i}\beta_{i}\,k \odot m_{i} + \tfrac12(\gamma_{11} + \gamma_{22})\,k \odot n ,
\]
the transverse terms cancelling, and each term on the right is a displacement ($k \odot \xi$ with $\xi = \alpha k$, $\beta_{i}m_{i}$, $\tfrac12(\gamma_{11} + \gamma_{22})n$). Displacements have zero curvature, so $\mathcal{R}_{k}[\mathcal{T} - \tfrac12\id\tr\mathcal{T} - P_{k}\mathcal{T}] = 0$ at null $k$. Now let $k = \omega e_{0} + \kappa$, $m_{1}, m_{2}$ orthogonal to $e_{0}$ and $\kappa$ (so orthogonal to $k$ for every $\omega$), and $\mathcal{T}(\omega)$ polynomial in $\omega$, conserved at $k$. Then $F(\omega) = \mathcal{R}_{k}[\mathcal{T} - \tfrac12\id\tr\mathcal{T} - P\mathcal{T}]$ is a polynomial in $\omega$ (the curvature map is quadratic in $k$) which vanishes at the two null points $\omega = \pm\sqrt{-Q(\kappa)}$, where $P\mathcal{T} = P_{k}\mathcal{T}$; so $F$ is divisible by $(\omega^{2} + Q(\kappa)) = Q(k)$, and, the curvature map being linear in $h$, $\mathcal{R}_{k}[h] = \lambda\,\mathcal{R}_{k}[P\mathcal{T}]/Q(k) + \lambda\,F/Q(k)$ with $F/Q(k)$ a polynomial.

\emph{(iv).} In the expansion of (iii), $\ip{\mathcal{T}_{1}}{\mathcal{T}_{2}} = \gamma_{11}\gamma'_{11} + \gamma_{22}\gamma'_{22} + 2\gamma_{12}\gamma'_{12}$ (every pairing that involves $k \odot k$ or $k \odot m_{i}$ with a map free of $n$ vanishes, $k$ being null and orthogonal to $m_{i}$), and the traces are $-(\gamma_{11} + \gamma_{22})$, $-(\gamma'_{11} + \gamma'_{22})$, so $B = \tfrac12(\gamma_{11} - \gamma_{22})(\gamma'_{11} - \gamma'_{22}) + 2\gamma_{12}\gamma'_{12}$, which is the stated form. It is positive semidefinite of rank two and vanishes exactly on the sources of (ii). The signatures at the sampled non-null momenta are computed; at a timelike momentum the conserved sources are the symmetric maps of the rest space, on which $B = \ip{\mathcal{T}_{1}^{\circ}}{\mathcal{T}_{2}^{\circ}} - \tfrac16\tr\tr$ has five positive directions and the trace as the negative one.

\emph{(v).} By (ii), $\ip{g}{\mathcal{G}_{k}[h]} = \ip{h}{\mathcal{G}_{k}[g]} = 0$ for every displacement $g$, so $\mathcal{G}_{k}[h](k) = 0$ for every $k$ and $h$. The components of $\mathcal{G}_{k}[h]$ are polynomials of degree at most two in $\omega$. From $\omega\,\mathcal{G}_{k}[h](e_{0}) + \mathcal{G}_{k}[h](\kappa) = 0$, the polynomial $b \cdot \mathcal{G}_{k}[h](\kappa)$ is divisible by $\omega$, so $b \cdot \mathcal{G}_{k}[h](e_{0})$ has degree at most one; and, the map being symmetric, $e_{0} \cdot \mathcal{G}_{k}[h](e_{0}) = -\kappa \cdot \mathcal{G}_{k}[h](e_{0})/\omega = \kappa \cdot \mathcal{G}_{k}[h](\kappa)/\omega^{2}$, a polynomial of degree at most two divided by $\omega^{2}$ and itself a polynomial: a constant. The degrees $0$, $1$, $2$ are attained (computed on the ten symmetric maps at $\kappa = 3e_{1} + 4e_{3}$).

\emph{(vi).} Computed on the six basis bivectors at three null momenta, of both handednesses and both signs of the frequency. The relation $\mathcal{R}_{k}[n](I_{4}B) = I_{4}\mathcal{R}_{k}[n](B)$ is the self-duality of a curvature with vanishing Ricci map. The comparison is with clause (2) of Proposition \claimref{field-theory.fock-space-maxwell-modes}: there the bivectors $f$ with $kf = 0$ are a real plane stable under $I_{4}$; here the curvature maps $\mathcal{R}_{k}[n]$, $n \in N_{k}$, are. The counts in $\Cl(1,2)$ (kernel of dimension three at two null momenta, equal to the displacements) and $\Cl(1,4)$ (kernel of dimension ten at one null momentum, the displacements and five transverse trace-free maps) are computed with the same formula for the curvature.

\emph{(vii).} For a conserved source at non-null $k$, $\theta_{k}$ is conserved, $\tr\theta_{k} = 3$, $\ip{\theta_{k}}{\theta_{k}} = 3$, and $\ip{\mathcal{T}^{\circ}}{\theta_{k}} = \tr\mathcal{T}^{\circ} - k \cdot \mathcal{T}^{\circ}(k)/Q = 0$, so $\ip{\mathcal{T}_{1}}{\mathcal{T}_{2}} = \ip{\mathcal{T}_{1}^{\circ}}{\mathcal{T}_{2}^{\circ}} + \tfrac13\tr\tr$ and $B = \ip{\mathcal{T}_{1}^{\circ}}{\mathcal{T}_{2}^{\circ}} - \tfrac16\tr\tr$. For $\mathcal{T}_{i}(a) = m_{i}(a \cdot e_{0})e_{0}$ at spatial $k$: $\mathcal{T}^{\circ}_{i} = m_{i}\bigl((a \cdot e_{0})e_{0} - \tfrac13\theta_{k}\bigr)$, $\ip{\mathcal{T}_{1}^{\circ}}{\mathcal{T}_{2}^{\circ}} = m_{1}m_{2}(1 - \tfrac23 + \tfrac13) = \tfrac23 m_{1}m_{2}$. The comparison bilinear $B_{N}$ is a definition. Its reading is argued, and no proved statement uses it: a field varied within the Weyl-flat class, $h = \phi\,\id$ modulo displacements, pairs with a source through its trace only, and its field equation is one condition on the trace of the source (clause (iii) of the stored first-order theorem); for a Lorentz-invariant equation of second order the first-order interaction is then a constant times $\tr\mathcal{T}_{1}\tr\mathcal{T}_{2}/Q$, which is the comparison bilinear of the statement; the constant that gives $-Gm_{1}m_{2}/r$ for bodies at rest, as in clause (ii) of the stored rate result (the static coupling of Proposition \claimref{field-theory.operator-newtonian-coupling}), is that of $B_{N} = \tfrac12\tr\tr$ with the same $8\pi G$. With $B_{N}$ the summed square for a pair of scalars of trace $(1 - 6\xi)s$ into another is $\tfrac14\tau_{i}\tau_{j}s^{4}$, independent of the angle; its integral over the sphere with the factor $\tfrac14$ for two pairs of identical quanta is $\tfrac{\pi}{4}\tau_{i}\tau_{j}s^{4}$, that is $\tfrac{\pi}{180}\cdot 45\,\tau_{i}\tau_{j}$, nine times the trace term of the stored weights (the amplitude $\tfrac12$ against $-\tfrac16$); for the trace-free sources of the stored result's clause (iii) it is zero. Inserting these weights in the stored result's clause (v): $\alpha = \tfrac{\pi^{3}}{17280}\cdot\tfrac{45\,\tau_{j}\sum\tau_{i}}{225}$, which for a scalar and four bath scalars is $\pi^{3}/21600$, and summed over the four scalars of a copy $\pi^{3}/5400$; the ratios to $5789\pi^{3}/82944000$ and $1519247\pi^{3}/82944000$ are $3840/5789$ and $15360/1519247$. The hypotheses are those of the stored clause (v), with (OS) read for a trace-only source.

\emph{(viii).} In Theorem \claimref{synthesis.dimmed-burst-time-rescaled-copy} the ratio is $D = \Psi_{\mathrm{tot}}/\Psi_{0}$ with $\Psi_{0}$ the coefficient of the standard expectation, and its clause (ii) holds for any such ratio, the visible emission being the same function of the mass. With the mirror copying this tower, $\Psi_{\mathrm{tot}} = \Psi_{\mathrm{vis}}(2 + g_{\mathrm{h}} + o)$ and $\Psi_{0} = \Psi_{\mathrm{vis}}(1 + g_{\mathrm{s}})$. The difference of the two options is $g/(1 + g)$. The value $g = 2/98$: of the power of a hole emitting only massless neutrinos, photons and gravitons, $2\%$ is in gravitons \cite{Page1976} (abstract, read on the web on 18 November 2026), so relative to the rest $g = 2/98$; then $(2 + g)/(1 + g) = 99/50$ and $2/(1 + g) = 49/25$. The fraction is for that mix of emitted species; with more species emitted $g$ is smaller, and for a hole emitting photons alone it is larger (not computed).

\emph{Reading.} That a rotation of $m_{1}, m_{2}$ rotates $n_{+}, n_{\times}$ by twice the angle is in the proof of (i). The statement that the deflection of light tests the trace-free channel is standard, cited from memory, and concerns results outside the store. The statements about emission from a ground state and about the order $G\,T^{2}$ of the emission excited by the matter's stress are standard and are cited from memory; they are labelled argued and no clause above uses them.
